\section{从合并起点到当前模型的修正需求}
\label{sec:motivation}

模型合并将同一预训练模型独立微调得到的专家整合为单一模型。权重平均与任务算术提供了直接的参数组合方式~\citep{wortsman2022model,ilharco2023editing}；Fisher merging、RegMean 及 MaTS 进一步利用任务相关的敏感度或输入统计，区分不同参数方向的作用~\citep{matena2022fisher,jin2023regmean,tam2024mats}。在已有合并结果上，我们关注一个后续问题：\textbf{如何识别当前模型仍然缺失的任务功能，并据此继续修正？}

一个促使我们关注这一问题的现象是：\textbf{起点的功能差距，可以在相同的后续修正流程中显著缩小。}在 CLIP ViT-B/32 八任务设置中，固定专家、无标签校准图像和修正规则，Average 与 RegMean 在同一独立诊断集上的准确率差距由 $22.36$ 缩小到四轮后的 $0.15$ 个百分点（图~\ref{fig:motivation}a）。这表明，常用合并起点留下的差异仍有相当部分可以被后续过程修复。我们因而将注意力放到\emph{修正过程中模型还需要什么}：合并起点提供初始能力，后续更新继续协调尚未满足的任务需求。以下为 seed 0 的配对机制诊断，每轮 CG100、最多四轮；完整协议见附录~\ref{app:protocol}。

\begin{figure}[!b]
  \centering
  \includegraphics[width=\linewidth]{motivation.pdf}
  \caption{\textbf{从起点适配到当前状态下的继续修正。}ViT-B/32，八任务，seed 0。
  (a) 三种常用起点在同一独立诊断集上的准确率。
  (b) Average 的完整测试准确率：刷新与冻结共享第一轮模型，均继续求解和选择步长；虚线为单轮 CG400。
  (c) 相邻学生状态的实际求解度量 $\bar A\otimes\bar G$ 的余弦距离，取层/任务中位数，纵轴为对数尺度。
  (a) 与 (b) 使用不同数据分割，数值不跨面板相减；(c) 比较同一批校准图像上的统计。完整对照见附录~\ref{app:protocol}。}
  \label{fig:motivation}
\end{figure}

进一步观察修正过程，Average 的完整测试准确率从第一轮的 $83.712\%$ 提高到第四轮的 $85.312\%$，说明一次接受更新后仍有继续修正的空间。但如何进行下一次修正同样重要：冻结初始统计后继续求解和选步，四轮得到 $84.519\%$；仅将初次求解延长到 CG400，得到 $84.565\%$（图~\ref{fig:motivation}b）。在这些预算下，继续求解初始统计问题没有复现刷新流程的全部收益。\textbf{这提示我们关注的不仅是求解走了多少步，还包括下一步依据哪个模型状态来测量。}

理解这一差异，需要考虑更新带来的网络状态变化。即使专家与图像固定，上游激活、学生预测和下游传播仍会变化。Average 首轮前后，固定诊断图像上的纠正梯度余弦中位数仅为 $0.248$；实际求解度量也发生明显变化（图~\ref{fig:motivation}c）。单看变化尚不能判断它是否影响更新，因此我们进一步固定\emph{同一个}第一轮 checkpoint，分别用当前与初始统计构造并接受修正：独立诊断 KL 的降低量为 $0.080560$ 与 $0.044534$。这一干预把统计的状态依赖与接受更新的实际效果联系起来，支持下一轮重新测量当前模型。

由此得到我们的设计原则：\textbf{依据当前模型的输入与师生分歧构造修正，接受后再适配新的网络状态。}对一层而言，输入决定参数改动作用于哪些特征，师生残差经当前下游网络反传后，指出仍有分歧的输出轴。我们分别用二阶矩 $A$ 和 $G$ 表达这两方面，以专家参数为锚点构造逐层合并问题，再用真实留出 KL 判断是否接受提案。这也明确了与专家统计的区别：专家处的测量服务于保留专家自身的行为，而我们希望衡量这些参数目标在\emph{当前合并网络}中的作用。因此，测量位置与接受后的重新测量是同一原则的两个环节。

上述观察支持的是有条件的适配：RegMean 上刷新与冻结的终态表现接近，后期统计也逐渐稳定；起点适配同样不意味着任意初始化都收敛到同一结果。附录~\ref{app:protocol} 保留这些边界，以及同状态干预的等幅对照和跨 block 干预，避免把接受终态的优势解释为所有原始更新方向都更优。
